\title{Thesis Proposal}
\author{
        Mark Reuter
}

\documentclass[12pt]{article}
\usepackage{mathrsfs}
\usepackage[utf8]{inputenc}
\usepackage[english]{babel}
\usepackage{amsthm}
\usepackage{amssymb}
\usepackage{amsmath}
\usepackage{tabularx}
\usepackage{ltxtable}

\newtheorem{theorem}{Theorem}[section]
\newtheorem{lemma}[theorem]{Lemma}

\theoremstyle{definition}
\newtheorem{definition}[theorem]{Definition}

\theoremstyle{remark}
\newtheorem*{remark}{Remark}

\usepackage[letterpaper, margin=1in]{geometry}
\usepackage{setspace}
\singlespacing

\usepackage{mathtools}
\DeclarePairedDelimiter\ceil{\lceil}{\rceil}
\DeclarePairedDelimiter\floor{\lfloor}{\rfloor}

\begin{document}
\maketitle

\singlespacing
\section{Introduction}
\doublespacing

Noise functions have been at the forefront of efficiently enriching the details of synthesized computer graphics. One such function is Ken Perlin's famous algorithm, Perlin Noise. Perlin Noise has seen widespread use in almost every aspect of procedural content generation\cite{lagaeetal2010}. Perlin Noise has several key advantages. It is fast, memory efficient, and capable of producing intricate patterns. But it is well known for having certain visual artifacts. One such artifact is its regular zero-crossing property. Perlin Noise will always evaluate to zero at integer lattice points. This artifact produces regularity in synthesized content. 

Two solutions to the zero-crossing property of Perlin Noise found in the literature include the techniques of Fractal Sums and Noise Distortion. Both of these techniques attempt to offset the lattice to avoid the zero-crossing property. But an analysis of the noise values produced by both techniques shows that they still produce highly regular noise. The noise values produced by these techniques express a limited range at the integer lattice points. While the noise values do not necessarily equal to zero, they do not produce a normal distribution around the noise.

A second popular noise function, Simplex Noise, also expresses a zero-crossing property. Simplex Noise has many of the same advantages of Perlin Noise. It is fast, memory efficient, and capable of producing intricate patters. Ken Perlin\cite{perlin2001} introduced Simplex Noise to remove some of the visual artifacts or Perlin Noise. But Simplex Noise still contains artifacts induced by its regular zero-crossing property. Unlike Perlin Noise which is based on the integer lattice, Simplex Noise is based on the simplex lattice. So Simplex Noise does not necessarily evaluate to zero at integer lattice points, but it will always evaluate to zero at simplex lattice points. 

I propose two alternatives to Perlin Noise and Simplex Noise. The first is based on the integer lattice, and the second is based on the simplex lattice. These functions maintain the memory efficiency and capability of producing intricate patterns of both Perlin Noise and Simplex Noise, but these proposed noise functions suffer from higher computational complexity. The advantage of these functions is that they express no regularity in their zero-crossings, and they express an ideal distribution of noise values at any point of evaluation.

The topic of this thesis is an investigation into the zero-crossing properties of Perlin Noise and Simplex Noise. This investigation is in two parts. First, I analytically demonstrate the consequences of the zero-crossing properties of Simplex Noise and Perlin Noise. Second, I give examples of procedural content which have visual artifacts induced by these zero-crossing properties. 

\singlespacing
\section{Related Works}
\doublespacing

\subsection{Noise Functions}

\paragraph{Perlin Noise}

Perlin Noise\cite{perlin1985} has been an ubiquitous noise function in compter graphics since its unveiling in 1985. Originally this function was designed for use in texture synthesis\cite{perlin1985, perlin1989}, but its applications to other areas of procedural generation have been well explored\cite{lagaeetal2010, smelik2014}. Perlin noise is a lattice gradient noise function. A \textit{lattice gradient noise function} generates noise by interpolating vectors, called \textit{gradients}, defined at the integer lattice points. Perlin's original interpolation function used the cubic polynomial $3t^2-2t^3$\cite{perlin1985}, but this function is discontinuous in the second derivative at integer points. For this reason, Perlin changed the interpolation function to the quintic polynomial $interp(t)=6t^5-15t^4+10t^3$, and this function has a zero first and second derivative at both $t=0$ and $t=1$\cite{perlin2002}.

Intuitively, Perlin Noise is computed as the weighted sum of $2^N$ gradients where $N$ is the dimension of the input vector $\textbf{v}$. The standard method for computing Perlin Noise defines an indexed set $G$ of gradients and a permutation $P$ of $\{1, ..., H\}$. $H$ can be any arbitrary positive integer, but implementations of Perlin Noise generally set $H=256$. The standard method for producing the weighted sum of gradients uses the $2^N$ integer lattice points surrounding $\textbf{v}$, named $\textbf{z}_1$, ..., $\textbf{z}_{2^N}$, where each $\textbf{z}_i = \langle z^i_1, ..., z^i_N \rangle$, and the gradients at the integer lattice points are obtained by computing the following equations:

\begin{align*}
	\textbf{g}_i = G[(P[z^i_1 \; \textnormal{mod} \; H] \oplus ... \oplus P[z^i_N \; \textnormal{mod} \; H]) \; \textnormal{mod} \; |G|] \\
	\textnormal{for each } 1 \leq i \leq 2^N, \\
	\textnormal{where } \oplus \textnormal{ is the bitwise exclusive-or.}
\end{align*}

These gradients $g_i$ are the ones on which Perlin Noise computes a weighted sum. The weights for use in this weighted sum are obtained using the interpolation function applied on the distances, named $\textbf{d}_1$, ..., $\textbf{d}_{2^N}$, between the integer lattice points to the input vector $\textbf{v}$. Specifically,the weights are the vectors $\textbf{c}_i$, that are computed as follows:

\begin{align*}
	\textbf{d}_i &= \textbf{v} - \textbf{z}_i\\
	\textbf{c}_i &= \textbf{d}_i * \prod_{j=1}^{N}(1-interp(|d^i_j|))
\end{align*}

Finally, Perlin Noise is computed as the sum of the dot product between the weight vectors and gradients:

\begin{equation*}
	perlin(\textbf{v}) = \sum_{i=1}^{2^N}(\textbf{c}_i \cdot \textbf{g}_i)
\end{equation*}

Perlin Noise has a regular zero-crossing property; if $\textbf{v}$ is identical to one of the integer lattice points, then $perlin(\textbf{v})=0$. The integer lattice points are not the only points at which Perlin Noise evaluates to zero; depending on definitions of $G$ and $P$, Perlin Noise can evaluate to zero elsewhere. Regardless of $G$ and $P$, though Perlin Noise always evaluate to zero at integer lattice points.

A noise function is \textit{band-limited} if it only produces frequencies within a band. Perlin Noise is not fully band-limited\cite{cookderose2005}. Patterns are constructed from \textit{bands} of noise which are limited to specific frequency ranges. Cook and DeRose have demonstrated that Perlin Noise produces both high and low frequencies outside of a band, and this leads to aliasing problems.

The choice of the gradients in the set $G$ can impact the quality of Perlin Noise in respect to the range of expressed noise values and band-limitedness. Originally,  Yoon et al.\cite{yoonetal2004} demonstrated that a 1-dimensional gradient set of two-thirds positive values and one-thirds negative values does not produce noise in the full range of $[-1, \; 1]$. Later, Yoon and Lee\cite{yoonlee2008} demonstrated that the choice of gradient set can produce more or less band-limited Perlin Noise.

\paragraph{Simplex Noise}

Ken Perlin\cite{perlin2001,gustavson2005} introduced another lattice gradient noise function called Simplex Noise. It is an improvement to his previous noise generation algorithm. Simplex Noise deviates from Perlin Noise in two key aspects. First, the integer lattice is replaced by the simplex lattice. Second, radial attenuation is used instead of polynomial interpolation on the gradients. These two modifications lower the computational complexity of the noise function in higher dimensions and remove the directional artifacts of Perlin Noise.

A \textit{$N$-simplex} is the simplest shape that tessellates $N$-dimensional space, and the \textit{simplex lattice} is the arrangement of these shapes to fill $N$-dimensional space. For 1-dimensional space, the simplex lattice is no different than the integer lattice. The simplex lattice in 2-dimensional space is a tessellation of equilateral triangles, and for 3-dimensional space, the simplex lattice is a tesselation of slightly skewed tetrahedra.

Simplex Noise is the weighted sum of $N+1$ gradients where $N$ is the dimension of the input vector $\textbf{v}=\langle v_1, ..., v_N\rangle$. The standard method for computing Simplex Noise defines an indexed set of gradients $G$ and a permutation $P$ of $\{1, ..., H\}$, and implementations of Simplex Noise typically set $H=256$. The standard method for producing weighted sums of gradients uses the $N+1$ simplex lattice points surrounding $\textbf{v}$, named $z_1, ..., z_{N+1}$, where each $z_i = \langle z^i_1, ..., z^i_{N+1}\rangle$. Computing the locations of the simplex lattice points surrounding $\textbf{v}$ is performed in two steps. Intuitively, the first step is to perform a skew operation on the simplex lattice which aligns the simplex lattice points with the integer lattice points. The effect of this skew operation is that $N$ factorial simplices form a $N$-dimensional, axis-aligned hypercube. The 2-dimensional simplex skew operation is defined as follows:

\begin{equation*}
	simplex2DSkew(\textbf{v}=\langle v_1, v_2 \rangle) = \floor{\textbf{v} + \frac{(v_1 + v_2) * (\sqrt{3} - 1)}{2}}
\end{equation*}

The 3-dimensional simplex skew operation is defined as follows:

\begin{equation*}
simplex3DSkew(\textbf{v}=\langle v_1, v_2, v_3 \rangle) = \floor{\textbf{v} + \frac{(v_1 + v_2 + v_3)}{3}}
\end{equation*}

A similar skew function can be defined for any arbitrary dimension. The second step is to decide which of the $N$ factorial simplices in which the input vector $\textbf{v}$ is contained(**work on this**). Where $\textbf{v'}=simplexNDSkew(\textbf{v})$, the ordering of the components of $\textbf{v'} - \floor{\textbf{v'}}$ correctly decides the skewed simplex containing $\textbf{v'}$. 

\begin{align*}
	\textbf{g}_i=G[(P[p^i_1  \; \textnormal{mod} \; H] \oplus ... \oplus P[p^i_{N}  \; \textnormal{mod} \; H]) \; \textnormal{mod} \; |G|] \\
	\textnormal{for each } 1 \leq i \leq N+1, \\
	\textnormal{where } \oplus \textnormal{ is the bitwise exclusive-or.}
\end{align*}

Further, Simplex Noise weights the contribution from the gradients using radial attenuation. Conceptually, radial attenuation is like a drop of water falling in a still pool. A ripple expands out from the drop, and it lessens in magnitude the further it gets from the drop. Radial attenuation weights the contributions of the gradients in a similar manner. The contribution of a gradient is strongest the closer the noise is sampled to it, and the contribution falls off to zero before the input vector crosses the boundary into a simplex further away. Formally stated, each $\textbf{p'}_i$ is the location of one of the $N+1$ vertices of the simplex containing the input vector $\textbf{v}$ and the two place function symbol, $e$, the Euclidean distance between two vectors squared. The constants $d_i$ are used to compute the falloff of the contribution from the gradients at the lattice locations. The weight vectors $\textbf{c}_i$ are defined as follows:

\begin{align*}
d_i &= max\{0, 0.5-e(\textbf{p'}_i,\textbf{v})\} \\
\textbf{c}_i &= d_i^4 * (\textbf{p'}_i - \textbf{v})
\end{align*}

Finally, the simplex noise function is computed as the sum of the dot product between the weight vectors and gradients along with a constant multiplication $C$ to yield noise in the correct range. 

\begin{equation*}
	simplex(\textbf{v}) = C * \sum_{i=1}^{N+1}(\textbf{c}_i \cdot \textbf{g}_i)
\end{equation*}

Like Perlin Noise, Simplex Noise has a regular zero-crossing property; if the input vector $\textbf{v}$ is identical to one of the simplex lattice points, then $simplex(\textbf{v})=0$. Also, these are not necessarily the only points at which Simplex Noise evaluates to zero; depending on the definitions of $G$ and $P$, Simplex Noise can evaluate to zero elsewhere. Regardless of $G$ and $P$, though Simplex Noise always evaluates to zero at integer lattice points.

To my knowledge, the impact of the choice of the gradient set $G$ using the techniques of Yoon and Lee has not been done in respect to Simplex Noise. As well, study of the band-limitedness of Simplex Noise using the techniques of Cook and DeRose is not complete.

\paragraph{Fractal Sums}

Fractal Sums are a method of producing fractal Brownian motion\cite{lagaeetal2010}. The Fractal Sum method computes the contribution of multiple octaves of noise. An \textit{octave} of noise is a pair $\langle f, a \rangle$ where $f$ is called the frequency and $a$ is called the amplitude. Implementations of Fractal Sums compute $F$ octaves of noise $\langle f_1, a_1 \rangle, ..., \langle f_F, a_F \rangle$, and each octave is related to the next by a lacunarity, $l$, and gain, $g$, as follows:

\begin{align*}
	f_1 &= 1, \\
	f_{i+1} &= f_i * l, \\
	a_1 &= 1, \\
	a_{i+1} &= a_i * g.
\end{align*}

In addition to these definitions of the octaves, typical implementations of Fractal Sums will define the lacunarity $l$ and gain $g$ in terms of each other. An implementation of Fractal Sums is called a $\frac{1}{l}$\textit{-noise} if $g=\frac{1}{l}$. Fractal Sums is not a noise function on its own. It relies on another noise function, such as Perlin Noise or Simplex Noise, as the basis. Fractal Sums samples another noise function at multiple frequencies and scales the noise value by the amplitude. Fractal Sums is defined formally with $F$ octaves of noise, $\langle f_1, a_1 \rangle, ..., \langle f_F, a_F \rangle$, and a one place noise function symbol $noise$ as follows:

\begin{equation*}
fractal(\textbf{v}) = \frac{\sum_{i=1}^{F}(a_i * noise(f_i * \textbf{v}))}{\sum_{i=1}^{F}a_i}.
\end{equation*}

This definition can inherit a regular zero-crossing property from the basis function. For example, if Perlin Noise is chosen as the basis function for Fractal Sums and an integer lacunarity is used, then the computed Fractal Sums will cross zero at each integer lattice point. Yoon and Lee\cite{yoonlee2008} investigated the use non-integer lacunarity when computing Fractal Sums of Perlin Noise. The use of non-integer lacunarity reduced the regularity of the zero-crossings of the fractal Brownian motion generated from Fractal Sums of Perlin Noise.

\paragraph{Noise Distortion}

Noise Distortion is a technique used by F. Kenton Musgrave\cite{musgraveetal1989} for two reason. First, it gives a different "feel" to the noise. Second, it helps over come the regular zero-crossing properties of Perlin Noise and Simplex Noise. Like Fractal Sums, Noise Distortion is not a noise function on its own. It relies on another noise function, such as Perlin Noise or Simplex Noise. Noise Distortion modifies the domain of a noise function. It does this by computing offsets for the input vector, then the technique evaluates the basis noise function at the offset input. Noise Distortion is defined formally with an input vector, named $\textbf{v}=\langle v_1, ..., v_N\rangle$, of dimension $N$, $N$ constants, named $c_1, ..., c_N$, and a one place noise function symbol, named $noise$, as follows:

\begin{equation*}
	noiseDistortion(\textbf{v})=noise(\textbf{v} + \langle noise(\textbf{v}+c_1), ..., noise(\textbf{v}+c_N) \rangle)
\end{equation*}

Noise Distortion reduces the regularity of the zero-crossing property of the basis noise function. If Perlin Noise is chosen as the basis function for Noise Distortion, then the Noise Distortion of Perlin Noise does not necessarily evaluate to zero at the integer lattice points. Noise Distortion has two distinct trade-offs. First, it is computationally expensive. To compute Noise Distortion, the basis noise function is computed $N+1$ times where $N$ is the dimension of the input vector. Second, Noise Distortion produces a noise that is visually distinct from Perlin Noise and Simplex Noise. The noise produced by Noise Distortion may be less ideal than Perlin Noise or Simplex Noise for some typical applications.

\singlespacing
\section{Proposed Work}
\doublespacing

\paragraph{Modified Perlin Noise} 

The generation of Perlin Noise only considers gradients selected from the $2^N$ integer lattice points surrounding the input vector. When the input vector falls on an integer lattice point, the spline based interpolation function causes the contribution from the other $2^N-1$ gradients to fall to zero. And the weight vector is based on the distance vector from the input vector to the integer lattice point. The contribution from the integer lattice point on which the input vector falls is also zero because the distance from the integer lattice point to the input vector is also zero. I propose a solution which considers more integer lattice points when generating noise and changes the interpolation function used.

This solution generates noised based on the $4^N$ integer lattice points surrounding the input vector. As well, radial attenuation is used instead of splines to interpolate the contribution from the gradients. This function is still defined as the sum of the dot product between weight vectors and gradients. The gradients $\textbf{g}_i$ are obtained in the same manner as Perlin Noise. The only difference being the number of gradients selected. Where $e$ is a two place function symbol denoting the Euclidean distance squared and $\textbf{z}_i$ are the surrounding integer lattice points, the weight vectors are redefined as follows:

\begin{align*}
d_i &= max\{0, 2 - e(\textbf{z}_i,\textbf{v})\} \\
\textbf{c}_i &= d_i^4 * (\textbf{z}_i - \textbf{v})
\end{align*}

And the proposed noise function is defined as follows:

\begin{equation*}
	noise_1(\textbf{v})=C*\sum_{i=0}^{4^N}(\textbf{c}_i \cdot \textbf{g}_i)
\end{equation*}

where $C$ is a constant multiple to obtain noise values in the range $-1$ to $1$. The zero-crossing property of Perlin Noise is caused by the contribution of all gradients to the noise value falling off to zero at integer lattice points. At integer lattice points, some of the gradients contributing to the noise value produced by $noise_1$ do necessarily fall off to zero. As in the definition of Perlin Noise, the contribution from the gradient at the integer lattice point on which the input vector falls is zero. But this noise function has a distinct advantage over Perlin Noise. At integer lattice points, not all gradients contributing to the noise value fall off to zero. So it is possible for the sum of the contributions to be non-zero. But the noise values produced at the integer lattice points tend to produce local minimums in the noise. This is due to the fact that the integer lattice point coinciding with the input vector has the largest contribution to the noise, but it necessarily falls to zero.

$noise_1$ calculates the contribution of many gradients whose weights have fallen off to zero. This is due to the fact that The contribution from the corners of the $2N$-hypercube intersect the innermost $N$-hypercube, but it does not intersect much of this hypercube. This leads to unnecessary computations. 

\paragraph{Modified Simplex Noise}

A downside to the proposed noise function $noise_1$ is its growth rate. It grows in $O(4^N)$ with the dimension $N$. Perlin Noise also has an exponential growth rate $O(2^N)$, but this growth rate is lessened by the fact that noise is generally evaluated in the second and third dimension and is rarely evaluated in dimensions higher than the fourth dimension. So while Perlin Noise does have exponential growth rate, it only computes the contribution from 4, 8, and 16 gradients in the second, third, and fourth dimensions respectively. So in the dimensions which Perlin Noise is most used, it is less computationally intensive than the growth rate might suggest. But the growth rate of $noise_1$ does pose a problem in these dimensions. $noise_1$ computes the contribution of 16, 64, and 256 gradients in the second, third, and fourth dimensions respectively.

One of the original reasons Ken Perlin introduced Simplex Noise was to reduce the growth rate of his original algorithm. To mirror this progress, I propose a second noise function the reduce the growth rate of $noise_1$ while still maintaining a reduced zero crossing property. This approach, like Simplex Noise, is based on the simplex lattice, and it maintains the use of radial attenuation as its interpolation function. There are two advantages to using the simplex lattice over the integer lattice. First, the contribution from fewer gradients is calculated. This noise function calculates the contribution from $(N+1)^2+(N+1)$ gradients rather than $4^N$ gradients. So the contribution from 12, 20, and 30 gradients are calculated in the second, third, and fourth dimensions respectively. Second, $noise_1$ calculates the contribution of many gradients which have fallen off to zero. This noise function calculates less of these zero contributions making the simplex lattice more favorable to the integer lattice.

The formal changes to Simplex Noise are minimal. The same skew function as simplex noise is used to convert the simplex lattice to an integer lattice for gradient selection. The only difference when selecting gradients is that, in addition to selecting gradients from the simplex contain the input vector, gradients are selected from the lattice points adjacent to those from the containing simplex. This amounts to the original $N+1$ gradients and an additional $(N+1)^2$ gradients. Let each $\textbf{p'}_i$ be one of these gradients, and let the two place function symbol $e$ be the Euclidean distance squared. The modification to the weight vectors is as follows:

\begin{align*}
d_i &= max\{0, 1.2 - e(\textbf{p'}_i,\textbf{v})\} \\
\textbf{c}_i &= d_i^4 * (\textbf{p'}_i - \textbf{v})
\end{align*}

This noise function is then defined as the sum of the dot products between these weight vectors and the selected gradients $\textbf{g}_i$

\begin{equation*}
noise_2(\textbf{v})=C*\sum_{i=0}^{(N+1)^2+(N+1)}(\textbf{c}_i \cdot \textbf{g}_i)
\end{equation*}

\paragraph{Analysis of Zero-Crossing}

I intend to demonstrate analytically the effects of the zero-crossing properties of Perlin Noise and Simplex Noise. The regular zero-crossings of these functions affect the entire noise function. When evaluated at points close to the lattice points, these noise functions do not express the entire range of noise values $[-1,\; 1]$. I propose tests of Perlin Noise and Simplex Noise which evaluates the distribution of the noise values at different locations in respect to the lattice points. I compare these results with the 

\paragraph{Artifacts Caused by Zero-Crossing}

In this work, I intend to demonstrate several visually significant artifacts caused by the zero-crossing properties of Perlin Noise and Simplex Noise. Then, I provide a comparison with the methods found in the literature for avoiding the zero-crossing property. Finally, I provide a comparison with my proposed noise functions.

The choice of texture basis function can heavily influence the procedural synthesis of textures. One feature that can be influenced is the periodicity of features the texture. The zero-crossing properties of Perlin Noise and Simplex Noise introduces visual periodicity. I propose the investigation into several common methods of procedural texture generation.

\singlespacing
\section{Work Schedule}
\doublespacing

\LTXtable{\textwidth}{schedule.tex}

\bibliographystyle{plain}
\bibliography{proposal}

\end{document}